\begin{definition}[Root number of a Dirichlet character] \label{definition:root_number_of_a_dirichlet_character}
\TODO{AI generated, verify}
Let $k \in \mathbb{Z}_{\geq 1}$, and let $\chi$ be a \CrefAndHyperrefIfExist{definition:dirichlet_character_modulo_a_nonnegative_integer}{Dirichlet character modulo $k$}. Let $a(\chi) \in \{0,1\}$ be the parity defined by $\chi(-1) = (-1)^{a(\chi)}$, i.e., $a(\chi) = 0$ if $\chi$ is even and $a(\chi) = 1$ if $\chi$ is odd. The \hldef{root number of $\chi$} is the complex number \hl{$\varepsilon(\chi)$} defined by
$$\hlin{\varepsilon(\chi) := \frac{\tau(\chi)}{i^{a(\chi)} \sqrt{k}},}$$
where $\tau(\chi)$ is the \CrefAndHyperrefIfExist{definition:gauss_sum_of_a_dirichlet_character}{Gauss sum of $\chi$}. If $\chi$ is \CrefAndHyperrefIfExist{definition:primitive_dirichlet_character}{primitive}, then $|\varepsilon(\chi)| = 1$.
\end{definition}